\documentclass[11pt]{article}

% ============================================================
% Packages
% ============================================================
\usepackage[utf8]{inputenc}
\usepackage[margin=1in]{geometry}
\usepackage{amsmath,amssymb,amsthm}
\usepackage{graphicx}
\usepackage{booktabs}
\usepackage{xcolor}
\usepackage{hyperref}
\usepackage{tikz}
\usepackage{enumitem}
\usepackage{pgfplots}

\pgfplotsset{compat=1.18}

% ============================================================
% Document Formatting
% ============================================================
\setlength{\parskip}{0.6em}
\setlength{\parindent}{0pt}

% ============================================================
% List Formatting
% ============================================================
\setlist[enumerate]{
    leftmargin=2em,
    itemsep=0.25em
}

% ============================================================
% Solution Environment
% ============================================================
\newenvironment{solution}
    {\par\noindent\textbf{Solution.}\par}
    {\par}

% ============================================================
% Link Styling
% ============================================================
\hypersetup{
    colorlinks=true,
    linkcolor=blue,
    urlcolor=cyan
}

% ============================================================
% Metadata
% ============================================================
\title{Ordinary Differential Equations: Abridged Notes}
\author{Chris Ritter}
\date{September 3, 2026}

% ============================================================
% Document
% ============================================================

\begin{document}
\maketitle
\begin{center}
\small\emph{These abridged notes were reconstructed afterward because I did not take
direct notes during class. They are not a record of the actual lecture.}
\end{center}
\tableofcontents
\newpage

\section{Topic Summary}
This reconstruction continues first-order modeling and introduces the
question of where solutions are guaranteed to exist. A model identifies
a changing quantity, balances its inflow and outflow, and then applies
an initial condition. For a linear equation, continuity of its
coefficients gives a unique solution on the interval containing the
initial point where those coefficients remain continuous. Nonlinear
models can require separate existence and uniqueness checks.

\section{Modeling with a Balance Law}
For an amount \(y(t)\) in a well-mixed tank,
\[
\boxed{\frac{dy}{dt}=\text{rate in}-\text{rate out}.}
\]
The outgoing concentration is the amount divided by the current
volume. If inflow and outflow rates agree, the volume stays constant.
This equation may be linear or nonlinear, depending on the rates.

\subsection{Example: Rinsing a Dye Tank}
\textbf{Adapted from Boyce et al., 11th ed., \S2.3, Problem 1.}
A 200-liter tank initially contains dye at 1 gram per liter. Fresh
water enters and mixed liquid leaves at 2 liters per minute. When is
the dye concentration 1\% of its initial value?
\begin{solution}
Let \(y(t)\) be grams of dye. The initial amount is 200 grams, and
the volume remains 200 liters. There is no incoming dye, while dye
leaves at rate \(2(y/200)=y/100\) grams per minute. Thus
\[
y'=-\frac{y}{100},\qquad y(0)=200.
\]
Separation gives \(y(t)=200e^{-t/100}\), so the concentration is
\(y/200=e^{-t/100}\) grams per liter. Setting it equal to 0.01 gives
\[
e^{-t/100}=0.01
\quad\Longrightarrow\quad
\boxed{t=100\ln100\approx460.5\text{ min}\approx7.68\text{ hr}.}
\]
\end{solution}

\section{Linear and Nonlinear Equations}
The first-order equation \(y'+p(t)y=g(t)\) is linear: \(y\) and its
derivative occur only to the first power and are not multiplied
together. Equations such as \(y'=y^2\) are nonlinear. A linear
initial-value problem has a unique solution on any interval containing
its initial point where \(p\) and \(g\) are continuous. This theorem
does not automatically extend across a zero of the coefficient of
\(y'\) before the equation is divided into standard form.

\subsection{Example: An Interval of Guaranteed Existence}
\textbf{Adapted from Boyce et al., 11th ed., \S2.4, Problem 1.}
For
\[
(t-3)y'+(\ln t)y=2t,\qquad y(1)=2,
\]
find an interval on which the linear theorem guarantees a unique
solution, without solving the equation.
\begin{solution}
For \(t\ne3\), the standard form is
\[
y'+\frac{\ln t}{t-3}y=\frac{2t}{t-3}.
\]
The logarithm requires \(t>0\), and both coefficients have a
singularity at \(t=3\). The largest interval containing \(t_0=1\)
on which both coefficients are continuous is
\[
\boxed{(0,3).}
\]
This is the interval guaranteed by the theorem; the calculation does
not claim that every solution must fail at its endpoints.
\end{solution}
\end{document}
